\documentclass[11pt]{article}
\usepackage[paperwidth=210mm,paperheight=297mm,textwidth=182mm,top=20mm,bottom=22mm]{geometry}
\usepackage{amsmath,amssymb,newtxmath}
\usepackage[no-math]{fontspec}
\setmainfont{Times New Roman}
\usepackage{xeCJK}
\setCJKmainfont[AutoFakeBold=2]{SimSun}
\usepackage{graphicx,booktabs,array,hyperref,placeins}
\hypersetup{hidelinks,unicode=true}
\raggedbottom
\setlength{\emergencystretch}{2em}
\renewcommand{\arraystretch}{1.15}
\usepackage{seqsplit}
\begin{document}
\title{Litz-Wire Modeling and Numerical Verification: Comprehensive Research Report · B Discretization and Postprocessing}
\author{Miao Yuxuan\\Huainan Normal University}
\date{}
\maketitle
\noindent 合并稿对应章节 / Chapter in merged manuscript: \href{report_en_single.pdf#chapter-B}{B}\par
\FloatBarrier
\section*{B Discretization and Postprocessing}
\hypertarget{chapter-B}{}
\pdfbookmark[0]{B Discretization and Postprocessing}{bookmark-chapter-B}
\subsection*{Triangle Derivation}
\hypertarget{B.triangle_derivation.heading}{}
\pdfbookmark[1]{Triangle Derivation}{bookmark-B.triangle_derivation.heading}
On a linear triangle, J=\(\sum\)\(_a\)NaJa. Current integration uses A/3 for each shape function. Joule integration uses the mass matrix and Joule expression in \href{report_en_single.pdf#E09}{(9)}. Material markers assign each cell to a strand, the return, or air. Coordinates are converted to metres before area evaluation, and the integrated power is independently compared with the archived scalar.


\subsection*{RT0 Face Flux and Lightweight Field Recovery}
\hypertarget{B.report_rt0.heading}{}
\pdfbookmark[1]{RT0 Face Flux and Lightweight Field Recovery}{bookmark-B.report_rt0.heading}
Original RT0 face fluxes are stored by direction and integer grid identity. For each cell axis, opposing fluxes q− and q+ give the center component (q−+q+)/(2h²) and local resistive power [|q−|²+|q+|²+Re(q−*q+)]/(3σh). Summing all three axes and all cells recovers total loss. Faces omitted at insulating sides have structural zero flux in the original discretization; this zero is not an independently solved boundary-flux validation.


Publication postprocessing independently reconstructs complex center J for all 27648 cells and compares it with the original cell\_fields.csv, while exact overlap integration checks archived loss. A section uses the actual discrete layer near the center and states its z coordinate rather than assuming that all samples lie exactly at 1 mm. The W/cm³ scale is obtained by multiplying SI loss density by 10⁻⁶. Vector sample identities and their common length scale are inspectable in the figure request and mapping data without an additional solver.


\hypertarget{E09}{}
\begin{equation}
\begin{gathered}
M_T=\frac{A_T}{12}\begin{bmatrix}2&1&1\\1&2&1\\1&1&2\end{bmatrix} \\
P_T=\mathbf J_T^H M_T\mathbf J_T/\sigma
\end{gathered}
\tag{9}\label{eq:E09}
\end{equation}
\end{document}
